% Brief (tikz-diagram-craft). Register row ch7-47 (must).
% Reader question: where does money sit at every step of a settlement, and
%   which single operation is allowed to create more of it?
% Claim: every operation redistributes value among exactly three buckets
%   (wallet, escrow, commons pool); only topUp changes the total supply S_P.
%   The worked partial-breach settlement in the text -- before (4,6,0), a
%   slash of p=2, after (8,0,2) -- instantiates this concretely.
% Evidence: website-v2/public/whitepaper/agent-transactions-whitepaper.tex
%   sec:conservation, Theorem (Conservation) and its proof sketch: topUp(u)
%   grows W_P by u; escrow(b) moves W_P to E_P; markRunning is a state
%   change inside E_P with no monetary movement; refund(id) moves E_P back
%   to W_P; slash(id,p) splits E_P into W_P (b-p) and C_P (p). The in-text
%   worked instance: before (W,E,C)=(4,6,0), slash p=2, after (8,0,2), both
%   summing to S_P=10 [verified].
% Must distinguish: the one edge that creates supply (topUp) from the four
%   that only redistribute it; the two-way split slash performs from the
%   single-target escrow and refund edges; the before/after instance from
%   the general grammar.
% Grammar chosen: typed stock-flow ledger (S12: buckets joined by pipes,
%   conserved sum set under the rails), operation badges keyed to an
%   aligned legend, with the worked numbers annotated under each bucket.
%   Rejected: a Sankey-width-proportional flow, which would need five
%   simultaneous widths the chapter never gives and would imply a
%   precision the model does not claim.
% Five-second test: "three buckets, one gold pipe creates money, four
%   plain pipes only move it, and 4+6+0 still equals 8+0+2."
\begin{figure}[H]
\centering
\pdfigurehue{pdgold}%
\begin{tikzpicture}[pd figure,x=1cm,y=1cm]
  \def\clwx{1.6}   % wallet x
  \def\clex{5.1}   % escrow x
  \def\clcx{8.6}   % commons pool x
  \tikzset{cl bucket/.style={pd bucket,minimum width=2.7cm,minimum height=1.3cm}}
  \node[cl bucket] (wallet) at (\clwx,0) {wallet\\$W_P$};
  \node[cl bucket] (escrow) at (\clex,0) {escrow\\$E_P$};
  \node[cl bucket] (commons) at (\clcx,0) {commons pool\\$C_P$};
  % topUp: the one edge that creates supply
  \draw[pd focus pipe] (\clwx-2.0,1.95) -- (wallet.north west);
  \node[pd focus label,anchor=south west] at (\clwx-2.05,2.05) {topUp$(u)$};
  \node[pd focus badge] at (\clwx-1.55,1.42) {1};
  % escrow(b): wallet -> escrow, top rail y=1.5
  \draw[pd pipe] ({\clwx+0.9},0.65) -- ({\clwx+0.9},1.5) -- ({\clex-0.9},1.5) -- ({\clex-0.9},0.65);
  \node[pd label,anchor=south] at ({(\clwx+\clex)/2},1.58) {escrow$(b)$};
  \node[pd badge] at ({(\clwx+\clex)/2},1.05) {2};
  % markRunning: self-loop above escrow only, no monetary movement
  \draw[pd hairline,densely dashed,->] (escrow.100) .. controls +(100:1.35) and +(80:1.35) .. (escrow.80);
  \node[pd note,anchor=south] at (\clex,2.28) {markRunning: no flow};
  \node[pd badge] at ({\clex+1.35},0.95) {3};
  % refund(id): escrow -> wallet, bottom rail y=-1.5
  \draw[pd pipe] ({\clex-0.9},-0.65) -- ({\clex-0.9},-1.5) -- ({\clwx+0.9},-1.5) -- ({\clwx+0.9},-0.65);
  \node[pd label,anchor=north] at ({(\clwx+\clex)/2},-1.58) {refund};
  \node[pd badge] at ({(\clwx+\clex)/2},-1.05) {4};
  % slash(id,p): escrow -> wallet, second bottom rail y=-2.4 (distinct
  % height from refund's); escrow -> commons, symmetric rail on the right
  \draw[pd pipe,draw=pderror,dash pattern=on 3.2pt off 2.2pt]
    ({\clex-0.5},-0.65) -- ({\clex-0.5},-2.4) -- ({\clwx+0.5},-2.4) -- ({\clwx+0.5},-0.65);
  \node[pd badge] at ({(\clwx+\clex)/2-.45},-2.80) {5};
  \node[pd breach label,anchor=west] at ({(\clwx+\clex)/2-.15},-2.80) {slash};
  \draw[pd pipe,draw=pderror,dash pattern=on 3.2pt off 2.2pt]
    ({\clex+0.5},-0.65) -- ({\clex+0.5},-1.5) -- ({\clcx-0.5},-1.5) -- ({\clcx-0.5},-0.65);
  \node[pd breach label,anchor=north] at ({(\clex+\clcx)/2},-1.58) {slash: $p\to C_P$};
  \node[pd badge] at ({(\clex+\clcx)/2},-1.05) {5};
  % worked instance, well clear of every rail
  \node[pd note,anchor=north] at (\clwx,-3.45) {$4 \to 8$};
  \node[pd note,anchor=north] at (\clex,-3.45) {$6 \to 0$};
  \node[pd note,anchor=north] at (\clcx,-3.45) {$0 \to 2$};
\end{tikzpicture}

\medskip
{\footnotesize\centering
$S_P := W_P + E_P + C_P$: before $4+6+0=10$, after $8+0+2=10$ --- unchanged.
\par}

\medskip
{\footnotesize
\begin{tabularx}{\linewidth}{@{}c l >{\raggedright\arraybackslash}X c@{}}
\toprule
 & operation & effect on the buckets & $\Delta S_P$ \\
\midrule
1 & topUp$(u)$ & $W_P \mathrel{+}= u$ & $+u$ \\
2 & escrow$(b)$ & $W_P \mathrel{-}= b$, $E_P \mathrel{+}= b$ & $0$ \\
3 & markRunning & state change inside $E_P$'s contributing set; no transfer & $0$ \\
4 & refund$(\mathit{id})$ & $E_P \mathrel{-}= b$, $W_P \mathrel{+}= b$ & $0$ \\
5 & slash$(\mathit{id},p)$ & $E_P \mathrel{-}= b$, $W_P \mathrel{+}= (b{-}p)$, $C_P \mathrel{+}= p$ & $0$ \\
\bottomrule
\end{tabularx}}
\caption{Only top-up increases supply. A partial breach moves $(W_P,E_P,C_P)$ from $(4,6,0)$ to $(8,0,2)$: the full escrow closes, while total supply stays 10. [internal, \protect\path{Conservation.tla}]}
\label{fig:bonded-conservation-ledger}
\end{figure}
